% !TEX root = ../main.tex
\section{Estratégia de Testes}
\label{sec:estrategia}

% ============================================================================
% 3.1 RESUMO
% ============================================================================

\subsection{Resumo da Estratégia de Testes}

O Educado verifica seu software por duas frentes que não se falam, e que também
não olham para o mesmo alvo. A frente automatizada mora dentro dos repositórios
de código \cite{educado-repos} e entra em ação quando uma mudança é proposta,
antes de ela ser aceita: verifica o código de toda a base, mas nunca exercita a
experiência de uso. A frente de campo é conduzida com quem usa a plataforma no
dia a dia, catadores de materiais recicláveis aprendendo educação financeira e
prevenção a golpe, e mede a experiência de uso de um único produto, o aplicativo
móvel, que é o que chega às mãos desse usuário \cite{psp5}; ela não toca no
código. A frente do repositório é estreita, restrita a unidades isoladas de
código. A de campo é mais ampla e mais formal, orientada a metas de qualidade
observadas diretamente com o usuário \cite{dashboard}. A separação entre as duas
é dupla, de método e de alvo, e por isso olhar só para os repositórios descreve
apenas parte do que o projeto verifica, e sobre um produto diferente do que o
catador usa. As subseções seguintes detalham níveis, tipos, automação,
organização, ambiente, regressão, integração contínua, revisão de código e
registro de defeito em cada uma dessas frentes; critérios de aprovação e
cobertura aparecem na Tabela~\ref{tab:estrategia-resumo}.

{\footnotesize
\begin{longtable}{>{\raggedright\arraybackslash}p{2.1cm} >{\raggedright\arraybackslash}p{4.3cm} >{\raggedright\arraybackslash}p{4.4cm} >{\raggedright\arraybackslash}p{1.6cm}}
\caption{Resumo da estratégia de testes identificada nas duas frentes.}
\label{tab:estrategia-resumo}\\
\toprule
\textbf{Elemento} & \textbf{O que foi identificado} & \textbf{Evidência} & \textbf{Avaliação} \\
\midrule
\endfirsthead

\toprule
\textbf{Elemento} & \textbf{O que foi identificado} & \textbf{Evidência} & \textbf{Avaliação} \\
\midrule
\endhead

Objetivos de teste
& Automatizada: exigência de teste em PR que muda comportamento; sem documento de plano ou política.
& \texttt{CONTRIBUTING.md} da api, seção ``Quality gate'' \cite{educado-repos}.
& Parcial
\\[0.1cm]

& Campo: objetivo SMART de ciclo de avaliação da qualidade, com prazo e público definidos.
& \cite[p.~10]{psp5}.
& Adequado
\\
\addlinespace
Níveis de teste
& Automatizada: somente unitário. A documentação declara testes de integração e de nível HTTP que não existem no código.
& 31 arquivos na api, 1 no web, 1 no app \cite{educado-repos}. Em \cite{educado-repos}: \texttt{getting-started.md}, seção Tests, e \texttt{tools.md}.
& Parcial
\\[0.1cm]

& Campo: sistema e aceitação, com validação da jornada completa junto ao usuário final.
& Roteiro do teste de usabilidade \cite[p.~17]{psp5}.
& Adequado
\\
\addlinespace
Tipos de teste
& Automatizada: funcional predominante; segurança presente mas não rotulada; sem desempenho, compatibilidade ou API.
& \cite{educado-repos}: 13 de serviço, 10 de validação, 4 de middleware, 1 de hashing de senha.
& Parcial
\\[0.1cm]

& Campo: usabilidade e regressão manual.
& \cite{ericsson1984,brooke1996,nielsen1993,garvin1998}.
& Adequado
\\
\addlinespace
Automação
& Automatizada: três stacks distintas, sem padronização nem justificativa registrada.
& \texttt{jest.config.ts}, \texttt{vitest.config.ts} e a chave \texttt{jest} do \texttt{package.json} \cite{educado-repos}.
& Parcial
\\[0.1cm]

& Campo: coleta manual em formulário, com instrumentação automática de sessão na segunda visita.
& Google Forms e Microsoft Clarity \cite[p.~16]{psp5}.
& Adequado
\\
\addlinespace
Organização da suíte
& Automatizada: arquivos junto ao código, sufixo \texttt{.test.ts}; convenção documentada só na api.
& \texttt{README.md} da api, seção ``Testes'' \cite{educado-repos}.
& Parcial
\\[0.1cm]

& Campo: roteiro por tela e checklist do entrevistador.
& Guia de entrevista e checklist \cite[p.~12]{psp5}.
& Adequado
\\
\addlinespace
Ambiente
& Automatizada: execução isolada, sem infraestrutura; dependências externas mockadas.
& \cite{educado-repos}: sem \texttt{setupFiles} no Jest, CI não sobe serviços, 16 dos 31 arquivos usam simulação de dependência.
& Parcial
\\[0.1cm]

& Campo: dispositivo Android real, na sede da cooperativa.
& Dispositivos registrados em \cite{dashboard}.
& Adequado
\\
\addlinespace
Regressão
& Automatizada: correção de defeito raramente acompanhada de teste.
& \cite{educado-repos}: 14 commits \texttt{fix}, dos quais 2 a 3 alteram teste.
& Ausente
\\[0.1cm]

& Campo: a segunda visita foi desenhada como reteste das melhorias da primeira.
& \cite[p.~10]{psp5}: 53 de 57 melhorias; 6 dos 15 indicadores não remedidos.
& Parcial
\\
\addlinespace
Integração contínua
& Automatizada: workflow nos três repositórios, disparado apenas por pull request.
& \texttt{ci.yml}, gatilho \texttt{pull\_request}, nos três repositórios \cite{educado-repos}.
& Parcial
\\[0.1cm]

& Campo: não se aplica.
& Não se aplica
& Não se aplica
\\
\addlinespace
Pull requests
& Automatizada: template exige declarar o teste realizado, sem verificação automática.
& Template de PR da api, seções ``How it was tested'' e checklist \cite{educado-repos}.
& Parcial
\\[0.1cm]

& Campo: não se aplica.
& Não se aplica
& Não se aplica
\\
\addlinespace
Code review
& Automatizada: uma aprovação exigida na api e no web; app sem proteção de branch.
& Proteção de branch via API do GitHub \cite{educado-repos}.
& Parcial
\\[0.1cm]

& Campo: não se aplica.
& Não se aplica
& Não se aplica
\\
\addlinespace
Defeitos e issues
& Automatizada: nenhuma issue nos três repositórios de código, contra 100 pull requests. Outros repositórios do projeto acumulam registros, concentrados em 2025.
& Listagem via API do GitHub, descontadas as pull requests, que a consulta inclui por padrão \cite{educado-repos}.
& Ausente nos repositórios de código
\\[0.1cm]

& Campo: 57 melhorias levantadas, priorizadas e rastreadas.
& Diagrama de Mudge \cite[p.~28]{psp5} e Balanced Scorecard \cite[p.~71]{psp5}.
& Adequado
\\
\addlinespace
Critérios de teste
& Automatizada: o CI não é verificação obrigatória, logo um PR com suíte falhando pode ser mesclado.
& Verificação de status não habilitada na api nem no web \cite{educado-repos}.
& Ausente
\\[0.1cm]

& Campo: meta numérica por indicador, entre 75\% e 95\%.
& 15 indicadores publicados em \cite{dashboard}.
& Adequado
\\
\addlinespace
Cobertura
& Automatizada: coletável na api, sem meta definida; ausente no web e no app.
& \cite{educado-repos}: script de cobertura, nenhum limiar configurado.
& Parcial
\\[0.1cm]

& Campo: cobertura de jornada, do cadastro à emissão do certificado.
& Roteiro do teste de usabilidade \cite[p.~17]{psp5}.
& Adequado
\\
\addlinespace
\bottomrule
\end{longtable}
}


% ============================================================================
% 3.2 OBJETIVOS DE TESTE
% ============================================================================

\subsection{Objetivos de Teste}

Nenhum dos três repositórios declara, em texto explícito, o que um teste deveria
alcançar. O objetivo existe, mas precisa ser reconstruído a partir de indícios
espalhados pelo fluxo de contribuição de código.

A frente de campo, em contraste, tem objetivo escrito com formato SMART:
``Realizar um ciclo de avaliação da qualidade do Educado, por meio de testes de
usabilidade nas frentes de Software, Conteúdo e Design, elaboração e análise de
indicadores de desempenho (KPIs), proposição e validação de melhorias,
envolvendo no mínimo 30 cooperados de 2 cooperativas até 30 de junho de 2026.''
Ele foi cumprido: as duas cooperativas previstas participaram, o mínimo de 30
cooperados foi superado e a segunda visita ocorreu em 22/06/2026 \cite[p.~10]{psp5}. É
o único lugar do projeto onde o objetivo de teste está no papel antes de o teste
acontecer.

\subsubsection*{Evidências}

O guia de contribuição do \texttt{educado-api} abre a seção ``Quality gate:
lint, tests, typecheck'' e, nela, afirma que qualquer pull
request que muda comportamento deveria vir acompanhada de teste. O
template de pull request cobra o mesmo por outro caminho: a seção ``How it was
tested'' pede que quem contribui descreva como testou, e o item
``Tests added or updated for the behavior I changed'' transforma isso
em caixa a marcar \cite{educado-repos}. É o rastro mais próximo de um objetivo
formal que o projeto tem.

\subsubsection*{Avaliação}

A documentação erra em dois sentidos distintos, e a distância entre erro e
realidade não é a mesma nos dois. O \texttt{educado-docs} promete teste de
integração e uso de \texttt{supertest} que não existem em nenhuma forma, nem
parcial: o exagero é total. O \texttt{educado-web} e o \texttt{educado-app}
afirmam não ter suíte de teste, e essa afirmação é incorreta na letra, porque
cada um tem um arquivo de teste e um CI que o roda a cada pull request, mas
próxima da verdade no conteúdo, já que esse arquivo confirma apenas que a
infraestrutura de teste funciona, não que o produto se comporta como deveria
\cite{educado-repos}. No \texttt{educado-app} a documentação, de 13/08/2026,
simplesmente antecede o arquivo de teste, criado em 18/08/2026. No
\texttt{educado-web} a documentação foi editada em 27/08/2026, nove dias depois
de o arquivo já existir, e manteve a afirmação: ou a informação ficou
desatualizada, ou a equipe considerou que um smoke test não constitui suíte. Só
o exagero do \texttt{educado-docs} engana de fato quem chega ao projeto sobre a
proteção que ele tem.

% ============================================================================
% 3.3 NÍVEIS DE TESTE
% ============================================================================

\subsection{Níveis de Teste}

A base é o único andar da pirâmide que existe em automação. Os três
repositórios de código têm teste unitário. O \texttt{educado-api}
\cite{educado-repos} concentra 31 arquivos, guardados ao lado do código que
verificam, cobrindo serviços de aplicação,
validações, middlewares HTTP, hashing de senha e cliente de armazenamento. O
\texttt{educado-web} \cite{educado-repos} e o \texttt{educado-app}
\cite{educado-repos} têm um arquivo cada, e os dois se declaram smoke test no
próprio cabeçalho: confirmam que a infraestrutura de teste liga, não que o
produto se comporta como deveria.

O nível de integração está vazio no código e prometido na documentação. Nenhum
dos três repositórios testa a costura entre peças, e o \texttt{supertest},
instalado no \texttt{educado-api}, nunca é chamado. O \texttt{educado-docs}
\cite{educado-repos} descreve o Jest fazendo unidade e integração, com
\texttt{supertest} no nível HTTP, e o handbook repete a promessa. Na prática, a
suíte roda com Postgres, Redis e S3 mockados, e o CI não sobe serviço externo
algum: se a aplicação parar de conversar direito com o banco, o cache ou o
armazenamento, isso só aparece quando alguém abre o aplicativo, não antes.

Sistema e aceitação existem, mas fora do repositório. As duas visitas do PSP5
\cite[p.~17]{psp5} validaram a jornada completa em Android real, contra metas entre
75\% e 95\%. É o mesmo caminho que um teste de ponta a ponta automatizado
percorreria, sem a repetibilidade de máquina: as ferramentas usuais do
ecossistema para isso não aparecem em nenhum dos três repositórios.

\begin{table}[H]
\centering
\caption{Níveis de teste identificados.}
\label{tab:niveis}
{\footnotesize
\begin{tabularx}{\textwidth}{>{\raggedright\arraybackslash}p{2.2cm} >{\raggedright\arraybackslash}X >{\raggedright\arraybackslash}p{3.0cm}}
\toprule
\textbf{Nível} & \textbf{Evidência} & \textbf{Avaliação} \\
\midrule

Unitário
& 31 arquivos em \texttt{src/**/\_\_tests\_\_} \cite{educado-repos}; um arquivo em cada front-end \cite{educado-repos}.
& Consolidado na api, incipiente no web e no app.
\\
\addlinespace

Integração
& Nenhum teste. O \texttt{supertest} está instalado e não é usado \cite{educado-repos}; a documentação declara o contrário \cite{educado-repos}.
& Ausente, e descrito como existente na documentação.
\\
\addlinespace

Sistema
& Jornada completa validada em campo, em dispositivo Android real \cite[p.~17]{psp5}.
& Presente apenas na frente manual.
\\
\addlinespace

Aceitação
& Validação com o usuário final contra metas por indicador, de 75\% a 95\% \cite{dashboard}.
& Presente apenas na frente manual.
\\
\addlinespace

End-to-end
& Nenhuma ferramenta de automação nos três repositórios.
& Ausente nas duas frentes.
\\
\addlinespace

Outro
& Smoke test declarado nos dois front-ends \cite{educado-repos}; teste de usabilidade formal em campo \cite[p.~14]{psp5}.
& Presente.
\\

\bottomrule
\end{tabularx}
}
\end{table}


% ============================================================================
% 3.4 TIPOS DE TESTE
% ============================================================================

\subsection{Tipos de Teste}

O funcional é o único tipo que o código exercita de fato. Dos 31 arquivos de teste do
\texttt{educado-api}, 13 cobrem serviço de aplicação e 10 cobrem validação de
entrada e de regra \cite{educado-repos}, sempre no mesmo molde:
entrada conhecida, resultado esperado. Falta o critério estrutural que decidiria
quando parar de testar. O projeto delimita onde a cobertura é coletada, restringindo a medição às
camadas de aplicação, de interface e de segurança da infraestrutura, mas não
define limiar mínimo de comando,
decisão ou caminho. A cobertura existe como número; não existe como régua.

Desempenho e integração com a API não são testados em nenhum lugar do código.
Nenhuma ferramenta usual de carga ou benchmark está instalada, e o
\texttt{supertest}, presente no
\texttt{educado-api}, nunca é chamado. O que existe é observação de campo com
critério objetivo, não opinião: o indicador de Desempenho do Sistema mede se a
resposta ocorre em até três segundos, e ficou em 68,8\% contra meta de 95\%
\cite{dashboard}, o que significa que em cerca de três de cada dez observações a
resposta ultrapassou esse tempo. É medição de comportamento real sob uso, não
teste de carga controlado. Compatibilidade segue o mesmo padrão de ausência:
nenhuma matriz de dispositivo definida em código, e a variedade de aparelhos
registrada em campo corresponde aos que foram usados nas sessões, não a uma
matriz de compatibilidade definida pelo projeto.

A segurança concentra o achado mais importante da subseção. No código, os
middlewares que cuidam de autenticação por token, autorização por papel,
exigência de conexão cifrada e rastreamento de requisição, mais a função de
hashing de senha, têm teste próprio
\cite{educado-repos}, sem que o projeto os nomeie como segurança. Fora do código,
o Dependabot está ativo nos três repositórios e já gerou 38 commits corrigindo
dependência com vulnerabilidade conhecida \cite{educado-repos}. Como nenhum
arquivo de configuração dessa varredura está versionado, ela foi habilitada nas
definições do serviço de hospedagem: é a única camada de segurança automatizada
dos repositórios analisados, e é invisível para quem só lê o repositório.

A acessibilidade e a clareza de interface expõem o contraste mais forte da
frente de campo. O indicador de acessibilidade mede uso autônomo, não visão ou
destreza, e ficou em 56,0\% contra meta de 85\% \cite{dashboard}: quase metade de
quem testou precisou de ajuda para operar a interface. A clareza de botões ficou
ainda mais distante da meta, 45,7\% contra 90\%. Esses são critérios de
comportamento observado, cumpridos ou não; a satisfação de quem usa o aplicativo
é medida à parte, por outro indicador do mesmo instrumento de campo, o que separa
duas perguntas fáceis de confundir: gostar do aplicativo e conseguir operá-lo
sozinho. Usabilidade é o tipo mais desenvolvido do projeto, com método formal em
campo, apoiado em protocolo verbal, escala de usabilidade, heurísticas de
interface e dimensões de qualidade
\cite{ericsson1984,brooke1996,nielsen1993,garvin1998}. Regressão existe e é
tratada
na subseção~\ref{sub:regressao}.

% ============================================================================
% 3.5 AUTOMAÇÃO
% ============================================================================

\subsection{Automação de Testes}

Cada repositório do Educado resolveu a automação de testes à sua maneira, sem
registro da escolha em ata ou guia de contribuição. A API usa Jest com o
adaptador ts-jest, ambiente Node e limpeza automática de mocks entre casos
\cite{educado-repos}. O repositório web usa Vitest com jsdom e a Testing
Library, com preparação carregada antes da suíte. O aplicativo usa o preset de
teste do Expo com a Testing Library para React Native, declarado no manifesto do
pacote. Entre os três, só a API já oferecia comando de coleta de cobertura; web e
aplicativo ofereciam apenas execução única. Para viabilizar a medição nos outros
dois, a equipe deste relatório adaptou ambos a produzir dados de cobertura,
apresentados na seção~\ref{sec:execucao}; a adaptação é intervenção da equipe,
não prática do projeto. Os três disparam a suíte no mesmo gatilho, a proposta de
mudança, ao lado de verificação de tipo e, em dois casos, de análise de estilo,
sempre antes de a mudança entrar.

Depois que a mudança entra, o aplicativo carrega outra automação: a
instrumentação da segunda visita roda em seu interior e captura tempo, sequência
de telas e toques sem intervenção humana \cite{dashboard}, não ferramenta externa
do pesquisador. Ela alimenta as duas frentes: a coleta de respostas em campo
segue manual, um pesquisador registrando em formulário enquanto outro conduz a
tarefa \cite[p.~16]{psp5}. Uma frente automatiza a verificação do código antes de
a mudança entrar, a outra automatiza a observação do uso depois, e as duas
convivem sem se falarem.

% ============================================================================
% 3.6 ORGANIZAÇÃO DA SUÍTE
% ============================================================================

\subsection{Organização da Suíte de Testes}

A API concentra a única suíte de tamanho relevante, e ali a organização segue um
padrão reconhecível: os testes ficam em pastas ao lado do código que cobrem, com
sufixo padronizado no nome do arquivo, agrupados por módulo de negócio, não por
tipo de teste \cite{educado-repos}. A distribuição por camada é desigual. A
camada de aplicação concentra 25 dos 31 arquivos; a de interface, 4; a de
infraestrutura, 2. Dos 19 módulos de negócio, 16 têm pasta de teste própria; três
não têm, catálogo, certificados e atividades do estudante.

A documentação da API declara que o teste mora ao lado do código correspondente.
Fora dela essa convenção não se sustenta. No repositório web, o único arquivo de
teste está numa pasta central dentro do código-fonte, enquanto o componente
exercitado mora em outra árvore de diretórios; no aplicativo repete-se o padrão,
com o arquivo de teste numa pasta na raiz e o componente correspondente em outro
ponto do projeto. Nenhum dos dois documenta convenção própria.

Simulação e fixture também não têm infraestrutura compartilhada nos três
repositórios: não há diretório de simulações, nem de dados de teste, nem de
fábricas de objeto. Na API, 16 dos 31 arquivos simulam dependência dentro do próprio arquivo, por substituição de módulo ou de método. Para o
volume atual isso não constitui falha, mas marca um teto: suíte que cresça sem
esse apoio tende a repetir preparação em cada arquivo.

O módulo de certificados figura entre os três sem teste, e a emissão do
certificado é o passo final da jornada validada em campo, com usabilidade
apontada entre os itens priorizados pela frente de campo \cite[p.~28]{psp5}. A
coincidência é digna de nota, sem que se afirme relação causal entre as duas
observações.

% ============================================================================
% 3.7 AMBIENTE
% ============================================================================

\subsection{Ambiente de Teste}

O ambiente que executa a suíte automatizada dos três repositórios do Educado não
depende de nenhum serviço externo. Nenhum deles exige banco de dados, cache,
armazenamento de objetos ou acesso de rede para rodar os testes, e nenhuma
variável de ambiente precisa ser fornecida de fora \cite{educado-repos}. A API
roda sobre o próprio interpretador Node; o web usa um navegador simulado; o
aplicativo usa o preset de teste do Expo, que fornece as simulações das APIs
nativas.

Essa ausência não vem de a aplicação ser simples. O produto usa Postgres, Redis e
um servidor de objetos compatível com S3 em funcionamento normal: a API traz um
arquivo de composição que sobe os três serviços para o ambiente de
desenvolvimento. Quem não precisa deles é o nível de teste que o projeto
automatiza. Como só existe teste unitário, cada peça é isolada e tudo em volta é
simulado, o que dispensa qualquer dependência real, inclusive com a biblioteca de
teste de nível HTTP instalada e sem uso. Se a suíte incluísse teste de
integração, a mesma infraestrutura que hoje só serve ao desenvolvimento teria de
subir também para testar; o ambiente de teste é simples porque o nível de teste é
raso, não porque o sistema dispense essas peças.

A versão da linguagem é fixada de três formas diferentes. A API e o web declaram
Node 20 num arquivo de versão que a integração contínua lê. O aplicativo, sem
esse arquivo, define o ambiente de desenvolvimento por um ambiente declarativo
Nix, e fixa a versão diretamente no fluxo de integração contínua.

O ambiente simulado não é neutro. O arquivo de preparação do web substitui a
implementação de armazenamento local do navegador simulado por uma versão em
memória, porque a implementação experimental do próprio Node pode sobrepor a do
navegador dependendo da versão instalada, quebrando o teste. A simulação
transfere o problema de ambiente em vez de eliminá-lo. Um teste de middleware da
API usa essa lógica a favor: altera o modo de execução entre desenvolvimento,
teste e produção para verificar a exigência de conexão cifrada em cada um, e
restaura o valor original ao final. Ali o ambiente é parâmetro controlado, não
dependência externa.

O teste de campo inverte cada um desses termos. Ocorreu na sede da cooperativa,
com um aparelho Android fornecido pela equipe de pesquisa e o aplicativo já
instalado \cite[p.~16]{psp5}. A segunda visita acrescentou instrumentação rodando
dentro do aplicativo, que registrou a sessão real \cite{dashboard}. É o único
momento em que o produto encontra um aparelho e uma pessoa reais.

% ============================================================================
% 3.8 REGRESSÃO
% ============================================================================

\subsection{Testes de Regressão}
\label{sub:regressao}

A prática de transformar defeito corrigido em teste permanente, mecanismo pelo
qual uma suíte de regressão cresce e deixa de repetir o mesmo erro, é exceção no
projeto Educado. Na API, 14 commits têm assunto iniciado por correção. Cinco
deles alteram algum arquivo de teste, mas dois tratam de ajuste de publicação sem
relação com lógica de aplicação, e um é commit de junção que duplica outro já
contabilizado. Sobram dois ou três casos em que a correção veio de fato
acompanhada de teste novo \cite{educado-repos}. O guia de contribuição da API
exige teste para toda mudança de comportamento; a prática observada, contudo, não
confirma essa exigência como rotina.

A automação de regressão existente nos três repositórios se resume à reexecução
da suíte a cada pull request. Não há execução agendada, disparo por
integração no ramo principal fora dessas propostas, nem verificação vinculada a
publicação de versão. Nenhuma versão publicada teve a suíte como portão, mas a
explicação é cronológica: as marcações de versão registradas nos três
repositórios datam de março de 2026, enquanto a integração contínua que executa a
suíte só foi criada em agosto. As publicações existentes antecedem a automação.

Quanto à manutenção, oito commits tocaram arquivos de teste na API, concentrados
em março, maio e junho, sem ocorrência posterior. O código de produção seguiu
curva equivalente: 17 commits em março, nenhum em abril, 2 em maio, 4 em junho,
nenhum em julho, 1 em agosto. O padrão indica repositório em
baixa atividade geral, não suíte abandonada enquanto o código evoluía
\cite{educado-repos}.

A frente de campo, por outro lado, executou o ciclo completo de regressão:
estabeleceu linha de base na primeira visita, acompanhou a implementação de 53
das 57 melhorias propostas e remediu os mesmos indicadores numa segunda visita
\cite[p.~10]{psp5}. Seis indicadores melhoraram e três pioraram, com 6 dos 15
medidos apenas uma vez, o que limita parcialmente a validação
\cite{dashboard}. Não há registro visível de regressão detectada pela suíte
automatizada; há registro de regressão real detectada por pessoas em campo.

% ============================================================================
% 3.9 CI/CD
% ============================================================================

\subsection{Integração Contínua}

Os três repositórios rodam fluxos automatizados a cada pull request, mas o
que cada um verifica difere. A API executa análise estática de estilo,
verificação de tipo e a suíte de testes. O aplicativo replica esse mesmo conjunto
de três etapas. Já o web substitui a análise de estilo pela construção do pacote
de distribuição, mantendo a verificação de tipo e a suíte de testes. Nenhum dos
três cobre exatamente o que os outros dois cobrem: só o web confirma que o
projeto ainda compila para publicação, e é o único que não audita o estilo do
código. Essa divergência não decorre de limitação técnica. Os três executam sobre
o mesmo tipo de máquina e a mesma versão da linguagem \cite{educado-repos}, o que
indica escolha de escopo por repositório, não restrição de ambiente.

O achado mais relevante está na obrigatoriedade do resultado. A API e o web
protegem o ramo principal e exigem aprovação humana antes de qualquer mesclagem,
mas nenhuma das duas proteções exige que as verificações automáticas tenham
passado: a configuração de proteção de ramo devolve, para os dois casos, que a
exigência de checagem de status está desabilitada \cite{educado-repos}. O
aplicativo não tem proteção de ramo alguma, apesar de acumular dezenas de
pull requests sob o mesmo fluxo. O guia de contribuição da API chama esse conjunto de verificações de
portão de qualidade \cite{educado-repos}; a configuração vigente implementa
aviso, não portão. Uma
proposta pode ser incorporada com a suíte falhando, desde que uma pessoa aprove.
A decisão final é sempre humana, sem poder de veto automatizado.

% ============================================================================
% 3.10 PR E CODE REVIEW
% ============================================================================

\subsection{Pull Requests e Code Review}

O modelo de pull request da API pede que quem contribui descreva como
testou a alteração e marque um item confirmando teste adicionado ou atualizado. O
guia de contribuição exige teste para toda mudança que altere comportamento, e a
API e o web protegem o ramo principal exigindo uma aprovação antes da mesclagem \cite{educado-repos}.

O histórico não confirma essa regra. Somando os três repositórios, 70 propostas
foram mescladas: 25 na API, 15 no web, 30 no aplicativo. Dessas, 38 partiram do robô de atualização de dependências e 32 foram de autoria
humana.
Nenhuma das 32 humanas tem revisão registrada, resultado checado por dois métodos
independentes, incluindo consulta individual a cada proposta. Das 38 do robô, 15
têm revisão registrada \cite{educado-repos}.

A configuração explica a lacuna sem contradizê-la. A proteção de ramo da API e do web exige
uma aprovação, mas a opção que estenderia essa exigência a quem tem privilégio
administrativo está desabilitada. Quem administra mescla sem
a aprovação que a regra pede, sem violar configuração nenhuma. O aplicativo não tem
proteção de ramo alguma \cite{educado-repos}.

O conjunto reúne os artefatos de um processo maduro de revisão sem a prática
correspondente nas mudanças que pessoas escrevem, e cobertura de teste ou
verificação automática também não são exigidas em nenhum ponto do fluxo. Revisão
de código é regra declarada e prática ausente; a configuração mostra como as duas
convivem sem contradição formal. Ausência de revisão registrada não prova
ausência de revisão, já que conversa fora da plataforma não deixa rastro. O que
se pode afirmar é que não há registro, e registro é o que torna a revisão
auditável.

% ============================================================================
% 3.11 ISSUES E DEFEITOS
% ============================================================================

\subsection{Defeitos e Issues}

Os três repositórios de código analisados, API, web e aplicativo, têm o mecanismo
de registro de defeito montado: modelo de relato de defeito, modelo de pedido de
funcionalidade, e um guia de contribuição que instrui quem contribui a abrir um
registro a partir desses modelos, verificando antes se o problema já não foi
relatado, para evitar duplicata \cite{educado-repos}. O mesmo guia orienta que
vulnerabilidade de segurança não seja aberta em registro público e aponta um
canal reservado, prática correta.

Nos três repositórios, porém, o número de registros abertos, em qualquer estado,
é zero. Essa ausência não reflete o histórico da organização. O repositório de
documentação, dentro do escopo desta análise, mantém 56 registros; um repositório
de aplicativo móvel de geração anterior tem 45; um repositório de backend
anterior, 17; um repositório de frontend anterior, hoje arquivado, 12
\cite{educado-repos}. No repositório de documentação a prática começa em março de
2025, concentra-se em setembro e outubro daquele ano, quando 46 registros foram
abertos, e depois cessa, com uma ocorrência em novembro e outra em dezembro.

O mecanismo foi transportado para os repositórios atuais; o hábito não. Isso não
é ausência de cultura de registro, é descontinuidade entre gerações do código. Na
API atual, 14 commits têm assunto iniciado por correção, evidência de que
defeitos existiram e foram corrigidos. O único registro deles é a mensagem do
commit, sem descrição do problema, sem passo de reprodução, sem vínculo com um
teste que impeça retorno.

O ciclo que falta no código sobrevive na frente de campo: 57 propostas levantadas
e priorizadas por comparação par a par, das quais 53 implementadas
\cite[p.~10]{psp5}, e o reteste da segunda visita detectando três indicadores que
pioraram \cite{dashboard}. No projeto, identificar, priorizar, corrigir e
verificar o efeito é ciclo completo só onde o código não é escrito.

